\section*{Реферат}
\pdfbookmark[0]{Реферат}{referat}

Отчёт содержит \ztotpages~стр., \total{figure}~рис. и \total{table}~табл. В списке литературы \total{reportreferences}~источников.

\noindent\textbf{Ключевые слова:} осадочный бассейн, теплоперенос, зрелость витринита, метод конечных объёмов, подвижная сетка, обратная задача, тепловой поток, нелинейный метод наименьших квадратов.

Рассмотрены прямая задача определения температурной истории и зрелости витринита и обратная задача восстановления нижнего теплового потока по современным скважинным наблюдениям. Работа выполнена в рамках постановки задачи Академического турнира 2026~года. Получены следующие результаты:
\begin{itemize}
\item сформулирована модель среды с изменяющейся геометрией, постепенным осадконакоплением и температурно-зависимыми свойствами;
\item построен численный алгоритм прямого расчёта и подбора истории потока;
\item разработана программная реализация с выводом полей на исходной геометрии и средствами ускорения повторных расчётов;
\item проведена верификация на задачах с известным решением и проверка восстановления потока по синтетическим наблюдениям;
\item выполнены расчёты трёх моделей, сопоставление с предоставленным эталонным решением и подбор потока по данным пяти скважин.
\end{itemize}

Математическая модель описывает теплопроводность, перенос тепла движущимся материалом и радиогенное тепловыделение. Пористость линейно меняется по мощности слоя и между конфигурациями. Начальное температурное поле определяется стационарной задачей. Зрелость вычисляется по модели EASY\%Ro с учётом момента отложения каждой материальной точки. В обратной задаче определяется пространственно постоянный, монотонно невозрастающий во времени нижний поток.

Численный алгоритм основан на методе конечных объёмов на подвижной колоночной сетке. Несовпадающие уровни ячеек соседних колонок учитываются многоточечной аппроксимацией потоков. Применяются неявный метод Эйлера, итерации Пикара и линейный метод GMRES с предобуславливанием ILU(2). Для подбора 17~параметров потока используется метод POUNDERS, работающий с отдельными компонентами вектора невязок без вычисления производных по параметрам.

Программная реализация разделяет расчётную часть на C++ и средства подготовки данных, вывода и визуализации на Python. Предусмотрены повторное использование геометрических данных, упрощение исходных границ и табулирование функции для вычисления интегралов зрелости. В трёх моделях выбранный итерационный метод сокращает время полного прямого расчёта в 1,31--2,55~раза относительно UMFPACK при рассмотренных параметрах и однопоточном выполнении.

Верификация методом изготовленных решений подтверждает порядок, близкий ко второму, на гладких стационарных задачах и ожидаемое уменьшение совместной пространственно-временной ошибки. На наклонном контакте материалов выявлено снижение порядка приблизительно до первого. В синтетической обратной задаче максимальные относительные невязки температуры и зрелости составляют 0,0029 и 0,0391\%, хотя отличие раннего потока достигает 11,01\%.

При сопоставлении Модели~1 с предоставленным эталонным решением относительные евклидовы нормы различий средних по слоям температуры и зрелости за все конфигурации равны 2,803 и 1,859\%. Для Модели~3 восстановлен поток по 500~температурам и 50~значениям зрелости в пяти скважинах. В контрольном расчёте в интервалы погрешности $\pm2\%$ попадают 456~температур и 39~значений зрелости. Согласование современных данных рассматривается отдельно от точности восстановления всей истории потока.
